\documentclass[8pt,a4paper]{article}
\usepackage[a4paper,landscape,margin=0.5cm]{geometry}
\usepackage[T1]{fontenc}
\usepackage[utf8]{inputenc}
\usepackage{lmodern}
\usepackage{amsmath,amssymb}
\usepackage{multicol}
\setlength{\parindent}{0pt}
\setlength{\parskip}{0.5pt}
\setlength{\baselineskip}{9pt}
\pagestyle{empty}
\renewcommand{\arraystretch}{0.8}

\newcommand{\E}{\mathbf{E}}
\newcommand{\B}{\mathbf{B}}
\newcommand{\D}{\mathbf{D}}
\newcommand{\Hf}{\mathbf{H}}
\newcommand{\J}{\mathbf{J}}
\newcommand{\A}{\mathbf{A}}

\begin{document}
\tiny

% PAGE 1: H2-H6
\begin{multicols*}{4}

\textbf{\footnotesize REM Exam Sheet -- Comprehensive Reference (H2-H6)}\\
{\tiny Constants: $c,\epsilon_0,\mu_0,k=1/(4\pi\epsilon_0)$}\\[2pt]


\textbf{H2 Electrostatica}\\
• 2.1 The electric field\\
• principle of superposition = interaction between any two charges is completely unaffected by the ...\\
• > calculate the forces on the charges separately, then the total force F = F 1+F2+...\\
• electrostatics = case where the source charges are stationary\\
• > the test charge may be moving\\
• 2.1.2 Coulombs law\\
• Coulombs law\\
• 2.1.3 the electric field\\
\columnbreak

\textbf{H3 Potentialen}\\
• charge above a conductor Consider a point charge q held a distance d above an infinite grounded c...\\
• > the potential above the plate is determined by q, but also the induced charge on the plate\\
• > this is a hard problem to solve, but we have the boundary conditions:\\
• instead, consider two point charges q at (0,0,d) and -q at (0,0,-d)\\
• > the potential is:\\
• > we can use the equation to solve our first problem\\
• 3.2.2 induced surface charge\\
• induced surface charge we can compute the surface charge via:\\
\columnbreak

\textbf{H4 E-velden in materie}\\
• H4: Electric fields in matter\\
• = contains 'unlimited' supply of charges and conduct electrical charge\\
• = all charges are attached to specific atoms or molecules\\
• > microscopic displacements aren't as dramatic, but does have a characteristic behaviour\\
• ionization place a neutral atom in an electric field\\
• polarizability tensor When the field is at some angle, we have to split into parallel and perpend...\\
• however, the most general linear relation between E and p is:\\
• force on polar molecules in\\
\columnbreak

\textbf{H5 Magnetostatica}\\
• 5.1 the Lorentz force law\\
• 5.1.2 magnetic forces\\
• Lorentz force law The magnetic force on a charge Q with velocity v in a magnetic field is:\\
• or thus in an E and B field:\\
• Work of magnetic force The magnetic force cant do work\\
• nl: if Q moves an amount dl = vdt, the work done is:\\
• current I = charge per unit time passing a given point\\
• > measured in Ampère: 1A = 1C/s\\
\columnbreak

\textbf{H6 B-velden in materie}\\
• H6: Magnetic fields in matter\\
• > apply a magnetic field B, the domains become aligned\\
• types of magnetization if we apply a magnetic field B, the we have 3 types of magnetization\\
• - ferromagnet: material retains magnetization after B is removed\\
• >> the magnetization isn't only dependent on the current applied field B\\
• > it is also dependent on the magnetic 'history' of the object\\
• 6.1.2 torques and forces on magnetic dipoles\\
• torque on a loop consider a rectangular current loop in a uniform field B\\
\columnbreak


\end{multicols*}

\vspace*{-10pt}
\newpage

% PAGE 2: H7-H12
\begin{multicols*}{4}

\textbf{\footnotesize REM Exam Sheet -- Comprehensive Reference (H7-H12)}\\[2pt]


\textbf{H7 Elektrodynamica}\\
• 7.1 electromotive force\\
• 7.1.1 Ohm's law\\
• conductivity \$\sigma\$  = ratio between current density J and force per unit charge f:\\
• resistivity \$\rho\$ = 1/\$\sigma\$\\
• Ohm's law The force f is due to the electromagnetic force, thus we can write:\\
• where vxB can be ignored, since the velocity is sufficiently small:\\
• Ohm's law in terms of V the current flowing from one electrode to another is given by:\\
• > R is the resistance with unit ohm $\Omega$\\
\columnbreak

\textbf{H8 Conservatiewetten}\\
• H8: Conservation laws\\
• 8.1 charge and energy\\
• 8.1.1 the continuity equation\\
• continuity equation The charge in a volume V is:\\
• and conservation of charge says:\\
• 8.1.2 Poynting's theorem\\
• Energy in E and B The work needed to assemble a static charge distribution is given by:\\
• And for the magnetic field:\\
\columnbreak

\textbf{H9 EM-golven}\\
• 9.2.1 the wave equation for E and B\\
• wave equation for E and B In regions of space with no charge or current, we have:\\
• which are two wave equations of waves travelling at the speed:\\
• with E0 and B0 the amplitude and \$\omega\$=kc the angular frequency\\
• > now Maxwell's equations put constraints on these waves:\\
• \$\nabla\$E=0 and \$\nabla\$B=0:\\
• \$\nabla\$xE=-\$\partial\$B/\$\partial\$t:\\
• $\nabla \times B = \mu_0 \epsilon_0 \frac{\partial E}{\partial t}$ doesn't imply any new constraints\\
\columnbreak

\textbf{H10 Potentialen en velden}\\
• H10: Potentials and fields\\
• scalar V and vector A potentials we can still define:\\
• however, last time for V we relied on \$\nabla\$xE, which is disproven due to:\\
• thus in a changing magnetic field E != -\$\nabla\$V\\
• thus, we can define:\\
• new Poisson's equation in a non-static case we thus have to account for the new potential formula...\\
• which reduces to Poisson's equation in static case\\
• we can put the relations for B and E in Maxwell's equation to find:\\
\columnbreak

\textbf{H11 Straling}\\
• 11.1.1 what is radiation?\\
• radiation when charges accelerate their fields can transport energy irreversibly out to infinity\\
• > imagine a giant sphere of radius r around this source\\
• > the power passing through its surface is the integral of the Poynting vector:\\
• > the em wave actually left the source an earlier time t0=t-r/c, thus:\\
• radiation of static sources The area of a sphere is 4\$\pi\$r\$^2\$\\
• > for radiation to occur the Poynting vector must decrease no faster than 1/r\$^2\$\\
• > according to Coulomb, electrostatic fields fall of like 1/r\$^2\$\\
\columnbreak

\textbf{H12 EM en relativiteit}\\
• principle of relativity = the laws of physics apply in any inertial reference frame\\
• 2: universal speed of light: speed of light in a vacuum is the same for all inertial observers\\
• Suppose A is on a train B going at a speed vBC with respect to the ground C\\
• > if i walk at a pace vAB with respect to the train, my speed with respect to the ground is:\\
• > for the observer on the train the time it takes for the light to hit the floor is:\\
• also define the quantity:\\
• > to an observer on the train, the time it takes is:\\
• is the time it takes to get to the front and \$\Delta\$t2 back again:\\
\columnbreak


\end{multicols*}

\end{document}
